\documentclass[conference]{IEEEtran}
\IEEEoverridecommandlockouts
\usepackage{cite}
\usepackage{amsmath,amssymb,amsfonts}
\usepackage{algorithmic}
\usepackage{graphicx}
\usepackage{textcomp}
\usepackage{xcolor}
\usepackage{subcaption}
\usepackage{float}
\usepackage{dblfloatfix}

\def\BibTeX{{\rm B\kern-.05em{\sc i\kern-.025em b}\kern-.08em
    T\kern-.1667em\lower.7ex\hbox{E}\kern-.125emX}}

\begin{document}

\title{Performance and Usability Evaluation of Continuous Hand Tracking via MediaPipe for Real-Time Character Control in Unity 2D}

\author{
  \IEEEauthorblockN{Pham The An}
  \IEEEauthorblockA{\textit{Information Technology Major} \\
    \textit{FPT University}\\
    anptse184524@fpt.edu.vn}
  \and
  \IEEEauthorblockN{Le Vu Truong Thinh}
  \IEEEauthorblockA{\textit{Information Technology Major} \\
    \textit{FPT University}\\
    thinhlvt...@fpt.edu.vn}
  \and
  \IEEEauthorblockN{Lam Nguyen}
  \IEEEauthorblockA{\textit{Information Technology Major} \\
    \textit{FPT University}\\
    lamn...@fpt.edu.vn}
  \and
  \IEEEauthorblockN{Tuan Dat}
  \IEEEauthorblockA{\textit{Information Technology Major} \\
    \textit{FPT University}\\
    dattn...@fpt.edu.vn}
  \and
  \IEEEauthorblockN{Le Do Quang Khoa}
  \IEEEauthorblockA{\textit{Information Technology Major} \\
    \textit{FPT University}\\
    khoaldq...@fpt.edu.vn}
}

\maketitle

\begin{abstract}
This paper evaluates both \textit{computational performance} and \textit{control usability} of continuous bimanual hand tracking in a Unity 2D arcade game using the homuler MediaPipe Unity Plugin (v0.16.3). Beyond FPS and CPU overhead, we measure gesture recognition accuracy, hand-detection reliability, and end-to-end response time from gesture onset to in-game action. Four runtime scenarios (keyboard baseline, single-hand, dual-hand, and stress) are tested on an ultrabook without discrete graphics. Results indicate interactive response with mean action latency of \textbf{XX.X}\,ms, overall gesture accuracy of \textbf{XX.X}\,\%, and an average FPS reduction of \textbf{XX.X}\,\% under dual-hand tracking relative to the keyboard baseline in Unity Editor Play Mode.
\end{abstract}

\begin{IEEEkeywords}
MediaPipe, Unity 2D, Hand Tracking, Gesture Accuracy, Response Time, Performance Evaluation, Real-Time Game Control.
\end{IEEEkeywords}

\section{Introduction}
Natural User Interfaces (NUI) enable camera-based game control without specialized hardware. While prior studies often report inference speed or frame rate, playable hand-controlled games also require \textit{accurate gesture mapping} and \textit{low action latency}. A system may run at acceptable FPS yet still feel unresponsive if landmarks are misclassified or actions arrive late.

Our project implements continuous bimanual control: left-hand finger counts drive movement and jumping; the right-hand index tip and grip gesture drive aiming and shooting. This paper therefore evaluates three complementary dimensions:
\begin{itemize}
\item \textbf{Performance:} FPS stability, CPU/GPU frame time, and pipeline latency.
\item \textbf{Usability -- gesture accuracy:} whether intended gestures produce the correct in-game command.
\item \textbf{Usability -- response time:} delay from gesture execution to observable game reaction.
\end{itemize}

Contributions: (1)~a deployed bimanual control framework; (2)~a four-scenario benchmark protocol; (3)~joint analysis of performance, gesture accuracy, and response time on consumer hardware.

\section{Related Work}
MediaPipe supports real-time hand landmark estimation~\cite{b1}, and the homuler plugin exposes Hand Landmarker to Unity~\cite{b2}. Game-vision studies often report detection accuracy or latency in isolation. We combine runtime performance with task-level gesture success rate and action latency inside an interactive 2D game, which better reflects player-perceived quality than FPS alone.

\section{System Architecture and Implementation}
The pipeline has three stages: webcam acquisition, native MediaPipe inference, and managed coordinate mapping with game input dispatch.

\subsection{System Overview and Data Pipeline Flow}
A webcam captures frames at 30\,fps. \texttt{WebCamSource} feeds \texttt{HandLandmarkerRunner}, which converts RGBA32 textures, applies flip/rotation, and invokes \texttt{HandLandmarker} (\texttt{numHands = 2}). \texttt{HandInputProvider} maps landmarks to \texttt{MoveX}, jump flags, aim position, and shoot state; \texttt{PlayerController} applies them to \texttt{Rigidbody2D} and \texttt{GunController}.

\IfFileExists{Image/Block-diagram.png}{%
\begin{figure}[htbp]
    \centering
    \includegraphics[width=\linewidth]{Image/Block-diagram.png}
    \caption{Flowchart of the real-time hand-tracking pipeline integrated into the Unity 2D game loop.}
    \label{fig:pipeline}
\end{figure}
}{}

\subsection{Multi-Hand Landmark Extraction}
Each hand yields 21 normalized landmarks $P_i=(x,y,z)$. Handedness tags distinguish left and right hands; labels are swapped when the webcam image is horizontally flipped.

\IfFileExists{Image/HandLandmark.png}{%
\begin{figure}[htbp]
    \centering
    \includegraphics[width=0.42\linewidth]{Image/HandLandmark.png}
    \caption{MediaPipe hand landmark topology. Landmarks 4, 5, and 8 support shooting and aiming logic.}
    \label{fig:landmarks}
\end{figure}
}{}

\subsection{Continuous Coordinate Mapping and Bimanual Control Logic}
\textbf{Screen mapping:} Landmark~8 maps to screen space via $X_{screen}=x\cdot W_{screen}$, $Y_{screen}=(1-y)\cdot H_{screen}$, then \texttt{Camera.ScreenToWorldPoint}.

\textbf{Bimanual logic:}
\begin{itemize}
\item \textbf{Left hand:} One finger $\rightarrow$ move left; two or more $\rightarrow$ move right; four fingers $\rightarrow$ jump.
\item \textbf{Right hand:} Landmark~8 defines aim; distance between Landmarks~4 and~5 below 0.06 triggers shooting.
\end{itemize}

These five commands---move left, move right, jump, aim update, and shoot---form the \textit{evaluation gesture set} used in Section~IV-C.

\section{Experimental Setup}

\subsection{Hardware and Software Environment}
Experiments were conducted on \texttt{DESKTOP-GNFSFFO}, an ultrabook-class PC with integrated graphics only; MediaPipe used the CPU delegate throughout (Table~\ref{tab:config}).

\begin{table}[htbp]
\caption{Experimental Hardware and Software Environment}
\centering
\begin{tabular}{|l|l|}
\hline
\textbf{Component} & \textbf{Specification/Version} \\
\hline
CPU & Intel Core i7-7600U @ 2.80 GHz \\
RAM & 12 GB DDR4 \\
GPU & Intel HD Graphics 620 (integrated) \\
Webcam & Built-in webcam, 720p @ 30 fps \\
Operating System & Windows 10 Pro, 64-bit (build 19045) \\
Game Engine & Unity 6000.4.6f1 \\
MediaPipe Plugin & homuler v0.16.3 \\
\hline
\end{tabular}
\label{tab:hardware}
\end{table}

\subsection{MediaPipe and Game Configuration}
\begin{table}[htbp]
\caption{MediaPipe and Runtime Configuration}
\centering
\begin{tabular}{|l|l|}
\hline
\textbf{Parameter} & \textbf{Value} \\
\hline
\texttt{numHands} & 2 \\
Running mode & \texttt{LIVE\_STREAM} \\
Image read mode & \texttt{CPUAsync} \\
Inference delegate & CPU (Windows) \\
Confidence thresholds & 0.5 / 0.5 / 0.5 \\
Target FPS / Camera & 60 / 30 fps \\
Resolution & $1920\times1080$ \\
\hline
\end{tabular}
\label{tab:config}
\end{table}

\subsection{Evaluation Scenarios}
\begin{itemize}
\item \textbf{S1 -- Baseline:} Keyboard/mouse; MediaPipe disabled (Fig.~\ref{fig:scenarios}a).
\item \textbf{S2 -- Single-hand:} \texttt{numHands=1}; move/jump gestures (Fig.~\ref{fig:scenarios}b--c).
\item \textbf{S3 -- Dual-hand:} Locomotion plus aim/shoot (Fig.~\ref{fig:scenarios}d--e).
\item \textbf{S4 -- Stress:} Dual-hand under combat or low-light conditions (Fig.~\ref{fig:scenarios}f--g).
\end{itemize}

\subsection{Performance Metrics}
Collected continuously during 5-minute sessions (three runs per scenario):
\begin{itemize}
\item \textbf{FPS:} average and minimum main-thread frame rate.
\item \textbf{Pipeline latency ($T_p$):} \texttt{OnResult} callback to valid output in \texttt{HandInputProvider}.
\item \textbf{CPU/GPU frame time:} Unity Profiler and Task Manager.
\end{itemize}

\subsection{Gesture Accuracy Metrics}
Accuracy was measured with a \textit{scripted gesture test} rather than free gameplay. For each gesture in Table~\ref{tab:gestures}, the tester performed $N=20$ repetitions under normal lighting. An observer recorded whether the expected game command occurred.

\begin{table}[htbp]
\caption{Evaluation Gesture Set and Expected Game Output}
\centering
\begin{tabular}{|l|l|l|}
\hline
\textbf{Gesture ID} & \textbf{Hand Pose} & \textbf{Expected Output} \\
\hline
G1 & Left: 1 finger & \texttt{MoveX} $\rightarrow -1$ \\
G2 & Left: 2 fingers & \texttt{MoveX} $\rightarrow +1$ \\
G3 & Left: 4 fingers & \texttt{JumpDown} = true \\
G4 & Right: index extended & Aim position updates \\
G5 & Right: thumb close to index & \texttt{ShootHeld} = true \\
\hline
\end{tabular}
\label{tab:gestures}
\end{table}

Three accuracy indicators are reported:
\begin{equation}
\text{GCA} = \frac{N_{\text{correct}}}{N_{\text{total}}} \times 100\%
\label{eq:gca}
\end{equation}
\begin{equation}
\text{HDR} = \frac{F_{\text{detected}}}{F_{\text{present}}} \times 100\%
\label{eq:hdr}
\end{equation}
where GCA is \textit{gesture command accuracy} and HDR is \textit{hand detection rate} when the hand is visible in frame. \textit{False activation rate} (FAR) is the percentage of idle frames that trigger an unintended command.

\subsection{Response Time Metrics}
Response time decomposes into three intervals (Fig.~\ref{fig:latency_breakdown}):
\begin{itemize}
\item \textbf{$T_1$ Inference:} gesture visible to camera $\rightarrow$ \texttt{OnResult}.
\item \textbf{$T_2$ Mapping:} \texttt{OnResult} $\rightarrow$ updated \texttt{HandInputProvider} output.
\item \textbf{$T_3$ Action:} updated input $\rightarrow$ observable game reaction (movement, jump animation, muzzle flash).
\end{itemize}
End-to-end action latency is $T_{\text{action}} = T_1 + T_2 + T_3$. For each gesture type, 20 samples were recorded using timestamp logs from \texttt{GestureEvaluationLogger}. A response was considered valid if $T_{\text{action}} \leq 150$\,ms, following common casual-game responsiveness guidelines.

\begin{figure}[htbp]
\centering
\fbox{\parbox{0.92\linewidth}{\centering\small
Gesture onset $\rightarrow$ [$T_1$] $\rightarrow$ MediaPipe result $\rightarrow$ [$T_2$] $\rightarrow$ HandInputProvider $\rightarrow$ [$T_3$] $\rightarrow$ Player action in game}}
\caption{Decomposition of end-to-end response time from gesture onset to in-game action.}
\label{fig:latency_breakdown}
\end{figure}

\begin{figure*}[htbp]
    \centering
    \begin{subfigure}[b]{0.24\textwidth}
        \includegraphics[width=\linewidth]{Image/KeyBoard.jpg}
        \caption{S1: Keyboard}
        \label{fig:scenarios:a}
    \end{subfigure}
    \hfill
    \begin{subfigure}[b]{0.24\textwidth}
        \includegraphics[width=\linewidth]{Image/OneHandMove.jpg}
        \caption{S2: 2-finger move}
        \label{fig:scenarios:b}
    \end{subfigure}
    \hfill
    \begin{subfigure}[b]{0.24\textwidth}
        \includegraphics[width=\linewidth]{Image/OneHandJump.jpg}
        \caption{S2: 4-finger jump}
        \label{fig:scenarios:c}
    \end{subfigure}
    \hfill
    \begin{subfigure}[b]{0.24\textwidth}
        \includegraphics[width=\linewidth]{Image/TwoHandMoveAndAim.jpg}
        \caption{S3: Move + aim}
        \label{fig:scenarios:d}
    \end{subfigure}
    \vspace{4pt}
    \begin{subfigure}[b]{0.24\textwidth}
        \includegraphics[width=\linewidth]{Image/TwoHandJumAndShoot.jpg}
        \caption{S3: Jump + shoot}
        \label{fig:scenarios:e}
    \end{subfigure}
    \hfill
    \begin{subfigure}[b]{0.24\textwidth}
        \includegraphics[width=\linewidth]{Image/RandomGesture1.jpg}
        \caption{S3: Dual-hand idle}
        \label{fig:scenarios:f}
    \end{subfigure}
    \hfill
    \begin{subfigure}[b]{0.24\textwidth}
        \includegraphics[width=\linewidth]{Image/RandomGesture2.jpg}
        \caption{S4: Combat stress}
        \label{fig:scenarios:g}
    \end{subfigure}
    \caption{Evaluation scenarios during Unity Editor Play Mode ($1920\times1080$).}
    \label{fig:scenarios}
\end{figure*}

\section{Performance and Usability Evaluation}

\subsection{Runtime Performance}
\begin{table}[htbp]
\caption{Runtime Performance (Average of 3 Runs)}
\centering
\begin{tabular}{|l|c|c|c|c|}
\hline
\textbf{Scenario} & \textbf{Avg FPS} & \textbf{Min FPS} & \textbf{$T_p$ (ms)} & \textbf{Drop (\%)} \\
\hline
S1 Baseline    & XX.X & XX.X & --   & 0.0 \\
S2 Single-hand & XX.X & XX.X & XX.X & XX.X \\
S3 Dual-hand   & XX.X & XX.X & XX.X & XX.X \\
S4 Stress      & XX.X & XX.X & XX.X & XX.X \\
\hline
\end{tabular}
\label{tab:results}
\end{table}

\subsection{Gesture Accuracy}
\begin{table}[htbp]
\caption{Gesture Command Accuracy (GCA) and Hand Detection Rate (HDR)}
\centering
\begin{tabular}{|l|c|c|c|}
\hline
\textbf{Gesture} & \textbf{S2 (\%)} & \textbf{S3 (\%)} & \textbf{S4 (\%)} \\
\hline
G1 Move left   & XX.X & XX.X & XX.X \\
G2 Move right  & XX.X & XX.X & XX.X \\
G3 Jump        & XX.X & XX.X & XX.X \\
G4 Aim         & --   & XX.X & XX.X \\
G5 Shoot       & --   & XX.X & XX.X \\
\hline
HDR (overall)  & XX.X & XX.X & XX.X \\
FAR (overall)  & XX.X & XX.X & XX.X \\
\hline
\end{tabular}
\label{tab:accuracy}
\end{table}

Overall GCA is the mean across G1--G5 (excluding inapplicable cells). Accuracy typically drops in S4 due to tracking loss under fast motion or poor lighting.

\subsection{Response Time}
\begin{table}[htbp]
\caption{End-to-End Response Time by Gesture (S3 Dual-Hand, ms)}
\centering
\begin{tabular}{|l|c|c|c|c|}
\hline
\textbf{Gesture} & \textbf{$T_1$} & \textbf{$T_2$} & \textbf{$T_3$} & \textbf{$T_{action}$} \\
\hline
G1 Move left   & XX.X & XX.X & XX.X & XX.X \\
G2 Move right  & XX.X & XX.X & XX.X & XX.X \\
G3 Jump        & XX.X & XX.X & XX.X & XX.X \\
G4 Aim         & XX.X & XX.X & XX.X & XX.X \\
G5 Shoot       & XX.X & XX.X & XX.X & XX.X \\
\hline
\textbf{Mean}  & XX.X & XX.X & XX.X & \textbf{XX.X} \\
\hline
\end{tabular}
\label{tab:response}
\end{table}

\subsection{Discussion}
\textbf{Performance vs.\ usability trade-off:} Lower FPS in S3--S4 does not automatically imply poor control quality. If $T_{\text{action}}$ remains below 150\,ms and GCA exceeds 85\%, the system remains playable for casual 2D games despite sub-60 FPS on ultrabook hardware.

\textbf{Error sources:} Misclassified finger counts (G1--G3) usually stem from partial occlusion or rapid transitions. Aim/shoot errors (G4--G5) increase when the right hand leaves the camera frame during jumping (Fig.~\ref{fig:scenarios}e).

\textbf{Feasibility:} Integrated GPU limits headroom; CPU-side \texttt{CPUAsync} dominates $T_1$. Kalman smoothing and higher-confidence thresholds may improve GCA at the cost of slightly higher $T_2$.

\section{Conclusion and Future Work}
We evaluated MediaPipe-based bimanual control along performance, gesture accuracy, and response-time axes. The framework is feasible on ultrabook hardware when mean action latency and GCA meet casual-game thresholds. Future work includes standalone builds, GPU-delegate testing, automated gesture labeling, and adaptive thresholds under low light.

\section*{Acknowledgment}
The authors thank faculty advisors at FPT University for their guidance.

\begin{thebibliography}{00}
\bibitem{b1} C.~Lugaresi \textit{et al.}, ``MediaPipe: A framework for building perception pipelines,'' arXiv:1912.01737, 2019.
\bibitem{b2} homuler, ``MediaPipe Unity Plugin v0.16.3,'' GitHub, 2024. [Online]. Available: https://github.com/homuler/MediaPipeUnityPlugin
\end{thebibliography}

\end{document}
